% ============================================================================
% SECTION 3 — SETUP: ENGINE, CORPORA, AND MATCHING FIBERS — BRIEF (contract)
% Sources: theory appendix Part II §0 (formal setup; the engine-grounded
%   fiber table), Part I §1 (M1 scaffold, kernels, rungs); prereg v2 §1.1
%   (truth asset), §4.1 (lineages/DGPs), §7.2 (enrichment); simulator
%   contracts doc for file:line anchors. Engine facts must match the frozen
%   artifacts; the bundle manifest is printed via \bundlemanifest{} (full
%   64-hex; never retype; Annex C-4 governs the known typo in companions).
% Must define, with macros from notation.tex: (1) the lab-asset-v3 engine —
%   exact integer uniform-price clearing, conservation + unit integrality +
%   tick lattice; engine kernels EXACTLY {fifo, random_unit_within_price};
%   pro_rata is NOT an engine kernel (S3 exact-arithmetic only, ruling
%   D1_14); 27 conformance tests + replay validator; read-only bundle.
%   (2) one-step scaffold: resting orders, pool \Pmarg, \Rstar, \Vstar,
%   rationed level, exhaustion boundary. (3) raw flow \rawflow and the
%   recorded map \recordedmap{}; the three grammar variants \Fexec,
%   \Fexecorders, \Ffull with the frozen statement that \Fexec{} IS the
%   corpus contract (under \Ffull{} the fiber is a point and T1–T3 are
%   vacuous — say so explicitly, prereg STOP-class reversal row). (4)
%   law-level vs realized-tape fibers \Fibkg{} vs \Fibtau{}; ambiguous set
%   \Utau; the per-kernel fiber-shape table (orthant / split subspace +
%   orthant / scale ray) from theory appendix Part II §0. (5) consumers
%   \Crisk, \Clat[W]; reference measure \muR{} (both preregistered
%   conventions: two-sided Unif[-R,R] split and one-sided box [0,R]
%   orthant). (6) lineages \Lone{} (ecomd_v2 transformer/potential; MACE
%   heritage — cite batatia2022 if invoked) and \Ltwo{} (recurrent
%   fact-surrogate, code-only pre-D0); DGPs D1 = lab-asset dual-arm + frozen
%   M0/M1-lumpable request generator, D2 = synthetic (Stage 2 only).
%   (7) enriched fixtures lab-asset-v3.1 (master seed 20260972; the frozen
%   21/22-event bundle has only single-maker touched levels — enrichment is
%   a hard prerequisite for S1a; never claim the bundle itself was mutated).
% FORBIDDEN: empirical results; describing \Ffull{} or \Fexecorders{} as the
%   training contract; any engine fact not verifiable in the grounding docs.
% ============================================================================
\section{Setup: engine, corpora, and matching fibers}
\label{sec:setup}

This section fixes the objects that every theorem package speaks about: the
frozen integer clearing engine and its two priority kernels
(\S\ref{sec:engine}--\ref{sec:scaffold}), the raw flow, the recorded map and
the frozen corpus contract (\S\ref{sec:corpora}), matching fibers
(\S\ref{sec:fibers}), consumers and reference measures (\S\ref{sec:consumers}),
flow-process classes (\S\ref{sec:policies}), and the lineages, data-generating
processes, and fixtures of the preregistered design (\S\ref{sec:lineages}).
Appendix~\ref{app:engine} carries the schema-level grounding (payload fields,
hash inventory, fixture facts); every engine fact used below is verifiable
there and in the frozen artifacts themselves. No statement in this section is
an empirical claim about any market: the engine is a laboratory object whose
semantics are frozen and hash-pinned.

\paragraph{Notation disambiguation.} The letters $K$, $D$, and $g$ each
serve several distinct frozen roles, fixed once here: $k$ is a priority
kernel; $\mathcal{K}_{\pi}$, $\mathcal{K}_1$, $\mathcal{K}_2$ and
$\hat{\mathcal{K}}$ are one-step aggregate-transition kernels of the T1
package and predictors thereof; $\Kdraw = 16$ is the frozen number of
inference-kernel draws of the design (\S\ref{sec:cube}); $P_K$ is the
projection onto a level's never-drawn split kernel and
$K_{\mathrm{split}}$ that subspace itself; $(N, K, n)$ are the
hypergeometric parameters local to Lemma~D (Appendix~\ref{app:proofs});
$D$ is a deployment map (T3); $\mathcal{D}$ is the set of price levels
deeper than the marginal level; $\mathsf{d}$ is a DGP instance (H5);
$\mathsf{Q}$ is a quotient of raw-flow space (V4); recording rungs always
carry a subscript ($\gagg$, $\gpo$, $\gpu$), and the letter $g$ without a
subscript is a policy response function (Definition~\ref{def:t1-piresp}).

\subsection{The lab-asset-v3 clearing engine}
\label{sec:engine}

The truth asset is the lab-asset-v3 bundle, frozen at its qualification exit
with manifest \bundlemanifest{} (canonical full 64-hex; a known companion-copy
typo is inventoried in the preregistration's errata register, Annex C-4). The
bundle is \emph{read-only}: it is never mutated, and every downstream artifact
--- including the fixture enrichment of \S\ref{sec:lineages} --- is a separate
object that cites it. Qualification comprised 27 conformance tests and a
deterministic replay validator that regenerates both serialized fixture arms
byte-exactly from disk.

The engine clears a single instrument on an integer tick lattice by exact
integer uniform-price matching. Volume, total cash, and total inventory are
conserved by every settlement transfer, and every fill is an integer number of
units. The books are order-level queues: a mapping from price to a list of
(order id, actor, remaining quantity) entries. Priority at a price level is one
of exactly two kernels --- the two members of the engine's
\fld{AllocationRule} enumeration:
\begin{itemize}
\item $\kfifo$ (engine name \fld{fifo}): arrival-ordered head-of-queue fill.
An incoming order consumes $\min(\text{residual demand},\ \text{head maker's
remaining quantity})$ from each successive pool order; the engine emits one
execution record per maker per incoming order and no draw payload.
\item $\kru$ (engine name \fld{random_unit_within_price}): unit-by-unit
clearing. Each executed unit is drawn uniformly without replacement from all
remaining resting units at the price: the engine computes the level's remaining
total, draws \fld{selected_unit} uniformly on that range, and walks the
per-order cumulative remaining quantities. Each unit emits one execution
record carrying \fld{maker_remaining} and the draw payload
$\fld{allocation_draw} = \{\fld{eligible\_units},\ \fld{maker\_order\_id},\
\fld{price},\ \fld{selected\_unit}\}$, with \fld{eligible\_units} the level's
remaining total before the draw.
\end{itemize}
$\kproprata$ is \emph{not} an engine kernel: pro-rata allocation appears in
this paper only as an exact-arithmetic taxonomy device (a frozen
preregistration ruling; internal ruling id anonymized for review) --- a
row of the static grid of \S\ref{sec:t1} and an arm of the S3
exact-arithmetic analysis (EP3, \S\ref{sec:cube-endpoints}), never an
executing arm.

Request validation is identity-dependent: self-trade prevention and per-actor
cash/inventory reserve checks reject requests
(\fld{self\_trade\_prevented}, \fld{insufficient\_cash},
\fld{insufficient\_inventory}). This channel is part of the engine's semantics
and enters \S\ref{sec:t1} through Proposition~R2-C.

The engine exposes two state hashes. The full state hash binds order identity,
queue order, and the RNG state. The \emph{aggregate statistic}
$A(s)$ --- the \fld{aggregate\_state\_hash} content --- records only the
per-price level sums on each side, total cash, total inventory, and the
last-match clock; by the schema's own contract, under an exogenous,
resource-slack request tape $A$ is equal after each completed request even
when identity allocations differ. $A$ is the strongest aggregate observer the
engine exposes, and the obstruction of \S\ref{sec:t1} is proved against $A$
itself; the same witness construction is visible in the level-total and print
projections (Appendix~\ref{app:t1proofs}).

\subsection{One-step clearing scaffold}
\label{sec:scaffold}

All theorem packages are instantiated on a one-step clearing round. Ask side
(with the bid side symmetric): $n$ resting orders with prices on the tick
lattice, quantities $q = (q_1,\dots,q_n) \in \mathbb{Z}_{\geq 0}^n$, arrival
clocks $t_1 < \dots < t_n$ within each price, and an actor map assigning each
order to an actor. An incoming marketable buy of integer size $Q$ walks the
book by price priority: every level whose cumulative supply stays below $Q$
clears fully; the \emph{marginal} (rationed) level $p^{*}$ carries the pool
$\Pmarg$ of its resting orders, pool size $\Rstar = \sum_{i \in \Pmarg} q_i$,
and residual demand $\Vstar = Q - S_{<} \in [0,\Rstar]$, where $S_{<}$ is
supply at strictly better levels. The \emph{interior regime} is the strict
partial fill $0 < \Vstar < \Rstar$; $\Vstar = \Rstar$ is the
\emph{exhaustion boundary} (the pool clears and the post-trade best price
jumps --- so the price-only post-quote reveals the stratum at every recording
rung). Levels deeper than $p^{*}$ (the set $\mathcal{D}$) are untouched.

The kernels allocate $\Vstar$ within $\Pmarg$ as follows. $\kfifo$
prefix-fills by arrival rank: with threshold index $\jstar$ the partially
filled pool order, the prefix allocation vector is $\cF$, and the executed
counts are a deterministic function of $(\qP, \Vstar)$. $\kru$ draws $\Vstar$
units sequentially, each uniform over the $\Rstar - t$ remaining pool units:
the per-step hit law is
$\mathbb{P}(\text{step } t \text{ hits order } i) =
(q_i - c_i^{(t)})/(\Rstar - t)$ with $c_i^{(t)}$ the hits so far, and the
per-order fill-count vector $c$ has the multivariate hypergeometric law
\begin{equation}
\label{eq:mvhg}
\mathbb{P}(c \mid \qP) \;=\; \MVHG(\qP,\Vstar)
\;=\; \frac{\prod_{i \in \Pmarg} \binom{q_i}{c_i}}
{\binom{\Rstar}{\Vstar}},
\qquad \textstyle\sum_i c_i = \Vstar ,
\end{equation}
the count marginal of the uniform-permutation sequence law (the ordering
given counts is ancillary). Deterministic kernels induce point-mass
allocations; \eqref{eq:mvhg} is the engine's only source of allocation
randomness.

The aggregate evolution of a round is kernel-independent: the level total at
$p^{*}$ decreases by exactly $\Vstar$, a level is deleted when its total
reaches zero, and settlement conserves totals. Everything the two kernels can
disagree about lives below the aggregate layer --- in the attribution of
$\Vstar$ to pool orders.

\subsection{Raw flow, recorded map, and the frozen corpus contract}
\label{sec:corpora}

The \emph{raw flow} of a round is $\rawflow = (\text{quantities},\
\text{arrival clocks},\ \text{actor map})$, an open subset of a
finite-dimensional space in the linear relaxation $\mathbb{R}_{\geq 0}$
(integrality is a lattice remark); a general linear parameterization
$(Q, q) = B\rawflow$ pulls every fiber back through $\ker B$
(rank--nullity: $\dim \ker B$ free directions are added; the pullback is
trivial only in the full-column-rank case, where $\ker B = \{0\}$ and the
parameterization changes nothing). The \emph{recorded map} $\recordedmap$ maps $\rawflow$ to
the tape $\tape$: the sequence of records the engine emits under kernel $k$
and recorded-field contract $\mathcal{F}$.

Three grammar variants of $\mathcal{F}$ were separated at design time, and the
choice is load-bearing:
\begin{itemize}
\item $\Ffull$: the entire event grammar, request events included. Under
$\Ffull$ the recorded map is injective on the raw flow (the requests
\emph{are} the flow), the fiber is a point, and Theorems T1--T3 are vacuous.
This is a preregistered STOP-class reversal row (the halt-with-no-repair
class of the preregistration; Section~\ref{sec:scope-stop}): re-admitting
request events into the corpus would not refute the theorems, it would
delete their subject.
\item $\Fexecorders$: $\Fexec$ plus the request events
(\fld{order\_request}/\fld{order\_accepted}). Since \fld{order\_accepted}
carries \fld{resting\_quantity}, every accepted order's declared quantity is
pinned and the fiber collapses to within-order timing ambiguity.
\item $\Fexec$: execution payloads plus $\fld{allocation\_draw}$ when present
plus pre/post best-bid/ask \emph{prices} (no sizes).
\end{itemize}
$\Fexec$ \emph{is} the frozen corpus contract (a frozen preregistration
decision; internal id anonymized for review): it is the
training corpus of every through-mechanism experiment in
\S\ref{sec:cube}, and every theorem in \S\ref{sec:t1}--\S\ref{sec:t3h5} is
stated under it. Within $\Fexec$, at each touched level with at least one
draw, the tape records: the draw sequence and drawn order ids;
\fld{maker\_remaining}, which pins each \emph{drawn} order's quantity exactly;
and \fld{eligible\_units} per draw, which pins the level's \emph{aggregate}
remaining quantity, hence the aggregate of the never-drawn orders --- the
engine records an aggregate where a per-order quantity would have removed the
symmetry. The fifo arm records no \fld{eligible\_units}; only executed orders'
\fld{maker\_remaining} is pinned. Untouched levels appear only through
best-price quotes: no depth at all.

The recording \emph{rungs} $g$ coarsen only the execution/allocation records:
$\gagg$ (aggregate prints, one per price: $\{(p, V_p)\}$; at $p^{*}$ the
kernel-independent statistic $\Vstar$), $\gpo$ (one fill record per filled
order: order id, price, $c_i$; zero-fill orders emit nothing), $\gpu$ (one
record per drawn unit, with the engine annotations). Ambient fields ---
pre/post BBO prices --- are held fixed across rungs, so the ladder is a pure
recording treatment. We write $T_{k,g}$ for the recorded map with execution
records further projected to rung $g$, and $\recordedmap$ for the unprojected
$\Fexec$ map.

\subsection{Matching fibers: law-level and realized-tape}
\label{sec:fibers}

Two identification objects must be distinguished throughout. The
\emph{law-level fiber} is
$\Fibkg(\rawflow_0) = \{\rawflow : \mathcal{L}_{k,g}(\rawflow) =
\mathcal{L}_{k,g}(\rawflow_0)\}$, where $\mathcal{L}_{k,g}$ is the
measure-valued map into laws of rung-$g$ tapes --- what a training
\emph{distribution} identifies. The \emph{realized-tape fiber} is
$\Fibtau = \{\rawflow' : \recordedmap(\rawflow') = \tape\}$ --- what one
round's record identifies. The \emph{ambiguous set} $\Utau$ is the set of
coordinates $\Fibtau$ leaves unpinned. The preregistered instrument explores
the law-level object engine-natively: a dual-hash fiber resampler regenerates
without-replacement interleavings holding the aggregate tape fixed
(Appendix~\ref{app:engine}). Table~\ref{tab:fibershapes} fixes the
engine-grounded fiber shapes used by T2 and T3.

\begin{table}[t]
\centering
\small
\begin{tabular}{@{}p{0.30\textwidth}p{0.31\textwidth}p{0.33\textwidth}@{}}
\toprule
Cell (kernel $\times$ tape) & Ambiguous set $\Utau$ & Fiber shape at
$\bar{\rawflow}$ \\
\midrule
$\kfifo$, $\Fexec$ & never-executed orders (any level) $+$ residuals of
partially executed orders beyond their last pinned \fld{maker\_remaining} &
translate of the orthant $\mathbb{R}_{\geq 0}^{|\Utau|}$ (recession cone) \\
$\kproprata$ (taxonomy only), aggregate & scale ray on the rationed side &
$\{c \cdot a : c \geq 1\}$ \\
$\kru$ $+$ per-unit tape, $\Fexec$ & at each touched level: split directions
among never-drawn orders $\{u : \sum_{i \in \mathrm{ND}(\ell)} u_i = 0\}$
(aggregate pinned by \fld{eligible\_units}; drawn orders pinned individually)
$\oplus$ never-executed orders at untouched levels & lower-dimensional
subspace (lineality) at touched levels $\oplus$ orthant at untouched levels
--- strictly smaller than the fifo fiber, never a point whenever $\geq 2$
never-drawn orders share a level or any level is untouched \\
$\kru$ $+$ aggregate tape & randomness marginalizes out; support union &
$\supseteq$ fifo fiber \\
\bottomrule
\end{tabular}
\caption{Engine-grounded fiber shapes at $\Fexec$ (theory appendix, Part II
\S0; $\mathrm{ND}(\ell)$ = never-drawn orders of level $\ell$). The pro-rata
row is taxonomy-only: $\kproprata$ never executes in this engine.}
\label{tab:fibershapes}
\end{table}

\subsection{Consumers and reference measures}
\label{sec:consumers}

A \emph{consumer} is a named market-native functional the deployed record is
asked to support. All consumers in this paper are finite linear functionals
$\Ccons(\rawflow) = v^{\top}\rawflow$ (assumption A3), with two named
instances: the raw-flow portfolio notional
$\Crisk(\rawflow) = \sum_i p_i q_i$ --- submitted/queue exposure, a risk
manager's object, \emph{not} the tape's cleared volume --- and the
latency-window notional $\Clat[W](\rawflow) = \sum_i p_i
\mathbf{1}\{\mathrm{clock}_i \leq W\}\, q_i$ --- what a truncating deployment
or a $\Delta t$-latency policy acts on. Allocation consumers
$C_w(a) = \sum_j w_j c_j$, linear in fill counts, are the deployment-side
readouts of \S\ref{sec:t3h5}.

The \emph{reference measure} $\muR$ puts product-uniform mass on the fiber
truncated to a box of radius $R$ in the ambiguous coordinates, with $R$ the
admissible excess scale (an energy-bounded adversary). Both per-coordinate
conventions are preregistered as frozen contract items and read as deviation
budgets: the \emph{one-sided box} on orthant coordinates,
$q_i \sim \mathrm{Unif}[a_i,\, a_i + R]$ independently ($q_i - a_i \in [0,R]$;
per-coordinate variance $R^2/12$, lower bound $a_i$ the recorded allocation or
pinned residual), and the \emph{two-sided split} convention on touched-level
split coordinates, $u = P_K \varepsilon$ with $\varepsilon_i \sim
\mathrm{Unif}[-R,R]$ i.i.d.\ ($|\varepsilon_i| \le R$; per-coordinate variance
$R^2/3$) and $P_K = I - \tfrac{1}{m}\mathbf{1}\mathbf{1}^{\top}$ the
projection onto the never-drawn split kernel of a level with $m$ never-drawn
orders. Gaussian references give the same scaling.

\subsection{Flow processes: the policy ladder}
\label{sec:policies}

The raw flow is generated by a stochastic process of actor requests. Two
classes matter. \emph{M0 (aggregate-lumpable)}: each actor's request at each
boundary is a measurable function of the anonymous aggregate history (level
sums, public information, own exogenous signals) plus exogenous policy
randomness; allocation identity never feeds back. \emph{M1$+$
(identity/resource-responsive)}: requests may additionally depend on the
identity-level components of the actor's own state --- own realized fills,
inventory, cash, induced values. (M2 feedback-adaptive and M3 rule-aware
policies are contained in M1$+$ for our purposes.) The boundary that the
theorems actually locate is stated in \S\ref{sec:t1}: ``no active identity
channel'', of which policy feedback (M1) and validation feedback (self-trade
prevention, per-actor reserves) are the two engine-native instances.

The data-generating processes of the preregistered design freeze the request
generator in the M0/M1-lumpable class: actions depend only on the anonymous
book state, own inventory, cash, and induced values --- never on allocation
identity --- so executed request streams are kernel-invariant and
request-level common random numbers pair the two arms exactly. This is the
single load-bearing pairing assumption of the experimental design; the
theorems of \S\ref{sec:t1}--\S\ref{sec:t3h5} do not rest on it (they are
statements about the engine's own laws), but every planned measurement does.

\subsection{Lineages, DGPs, and fixtures}
\label{sec:lineages}

Two simulator lineages carry the attribution cube of \S\ref{sec:cube}.
$\Lone$ is the transformer/potential surrogate lineage of
Section~\ref{sec:setup} (graph-potential architecture; internal codename
anonymized for double-blind review): state-dependent
interaction potentials in the graph-potential sense of MACE
\citep{batatia2022}, adapted to book state. $\Ltwo$ is a recurrent
fact-surrogate lineage --- code-frozen before D0, carrying no training claim.
Two data-generating processes are used: D1 = the lab-asset dual-arm engine
(same prestate and request stream, only \fld{allocation\_rule} exchanged)
driven by the frozen M0/M1-lumpable request generator; D2 = a synthetic
generator reserved for Stage 2. No outcome of any run is asserted anywhere in
this paper.

The frozen 21/22-event bundle instantiates the \emph{collapse stratum}: its
only touched level carries a single maker, so it contains no multi-order
rationed level. The enriched fixture family lab-asset-v3.1 (master seed
20260972; a separate artifact that cites, and never mutates, the frozen
bundle) exists precisely to place multi-order pools at touched levels on the
S3 scaffold $\qP = (5,4,3,2)$ with a $\Vstar$-ladder, plus clock-straddling
never-drawn sets; Appendix~\ref{app:engine} carries its contract. This
enrichment is a hard prerequisite for the post-D0 S1a gates (the
preregistered constructed-fiber-pair gate set, EP4 of
\S\ref{sec:cube-endpoints}) that would
dynamically confirm the multi-order cells; the static and law-level theorems
stand on the engine semantics alone. \RESULTSPENDING{}
